% Paper 9 of The Windstorm Series
\documentclass[12pt,a4paper]{article}

\usepackage{amsmath}
\usepackage{amssymb}
\usepackage{graphicx}
\usepackage{hyperref}
\usepackage{natbib}
\usepackage{booktabs}

\title{The Hardware Basin: Why the Quantization Cliff Is About Level Allocation, Not Bit Count}

\author{Grant Lavell Whitmer III\\
The Windstorm Institute, Fort Ann, NY 12827, USA\\
\texttt{grantwhitmer3@gmail.com}}

\date{April 2026}

\begin{document}

\maketitle

\begin{abstract}
Paper 7 reported a universal quantization cliff at INT4 to INT3 using bitsandbytes software quantization. We investigated whether this cliff is a software artifact or a mathematical property of low-precision arithmetic. The cliff location depends on the quantization method. Symmetric uniform quantization produces catastrophic BPT $\approx 17$ at INT4 (versus $4.3$ at FP16). bitsandbytes NF4 (normal-float-4, which places quantization levels at Gaussian quantiles) produces operational BPT $\approx 3.9$--$4.7$ at INT4. Same bit count, opposite outcomes. The difference is level allocation. A Lloyd-Max (minimum-MSE) quantizer achieves cosine similarity $0.990$ at INT4 and $0.965$ at INT3. The result is universal across four architectures and consistent across all 24 layers of Pythia-410M. Version 2.1 added three publication-grade verifications: five-seed replication at Pythia-1.4B (Welch $t = 633.74$, $p = 2.84 \times 10^{-15}$, Cohen's $d = 400.81$); Lloyd-Max INT3 end-to-end test showing that per-matrix cosine overstates propagated quality (BPT $= 11.74$); and three-model, three-seed robust replication. Version 2.2 (April 2026) adds the publication-readiness presentation: a methodological-journey narrative ($\S$1.4) tracing the seven rounds of follow-up --- including the Round-5 thesis pivot from ``cliff at INT4'' to ``cliff is about level allocation'' --- and an adversarial-review-defense table ($\S$4.6); Zenodo \texttt{10.5281/zenodo.19672922} (concept DOI \texttt{10.5281/zenodo.19672921}). The minimum viable inference specification is not an $N$-bit integer but an $N$-bit representation with distribution-aware level allocation, validated end-to-end.
\end{abstract}

\noindent\textbf{Keywords:} quantization cliff; level allocation; NF4; symmetric quantization; Lloyd-Max; inference hardware; structural bonus.

\section{Introduction}

The full text of this paper is provided as \texttt{paper.pdf} (current version v2.2) in this repository. The markdown source for the latest version is in \texttt{paper/Paper9-v2.2-draft.md} in this repository (with historical \texttt{paper/Paper9-v2.1-draft.md} alongside).

\section*{Code and Data Availability}

All code, experimental data, and figures are available under MIT license at \url{https://github.com/Windstorm-Institute/hardware-basin} and \url{https://github.com/Windstorm-Labs/hardware-basin}. The Grand Slam Supplementary Materials providing statistical verification (v2.1) and the v2.2 methodological-journey + adversarial-review-defense framework are at \texttt{grandslam\_supplementary.pdf} and \texttt{paper.pdf} (current version v2.2) in this repository.

\end{document}
